\section{Worked code for the literature catalogue}
\label{app:code}

This appendix gives runnable code for the assumption catalogue, the designs
in \cref{sec:illustrations}, the rows of \cref{tab:audit}, and further
instances from the literature: more treatment and outcome levels, two
treatments under one latent, and order-only chains. Each block is self-contained. Observational
tables come from \pid\texttt{.problems} or are written out; the graph,
query, and assumptions are explicit so that each block can serve as a
template for new data. The comment under each block is the printed output
(package version $0.1.0\mathrm{a}5$). Except for the last block, every
program below is certified LP or point by \cref{thm:lp}, and every interval
is \texttt{tight=True}.

\subsection{The assumption catalogue on the job-training table}
\label{app:code-zoo}

One table, one query, and every factory of \cref{tab:catalog} composed with
the causal order on~\eqref{eq:re1}. The order-only row is the natural bound
of \citet{manski1990nonparametric}; the other rows are the assumptions named
in that table.
\begin{lstlisting}
from partialid import Assumptions, Query, Variables, bound

V = Variables.from_binary("X", "Y")
P = {(0, 0): 0.30, (0, 1): 0.20, (1, 0): 0.15, (1, 1): 0.35}
Q = Query.interventional(outcome="Y", do={"X": 1}, value=1)
base = Assumptions.order("X", "Y")
catalogue = {
    "order (Manski)": base,
    "MTR": base & Assumptions.mtr("Y", "X"),
    "MTS": base & Assumptions.mts("X", "Y"),
    "MTR & MTS": base & Assumptions.mtr("Y", "X") & Assumptions.mts("X", "Y"),
    "Tan MSM, L=3": base & Assumptions.tan_msm("X", 3.0),
    "Rosenbaum, G=9": base & Assumptions.rosenbaum_gamma("X", 9.0),
    "c-dependence, c=0.1": base & Assumptions.c_dependence("X", 0.1),
    "margin, kappa=0.05": base & Assumptions.margin_budget(0.05),
    "f-sens. KL, rho=0.1": base & Assumptions.f_sensitivity("X", 0.1, f="kl"),
    "support k=3": base & Assumptions.support_bound(3),
}
for name, a in catalogue.items():
    r = bound(V, P, Q, a)
    out = f"[{r.lower:.3f}, {r.upper:.3f}]" if r.success else "infeasible"
    print(f"{name:22s} {out:16s} {r.method}")
# order (Manski)         [0.350, 0.850]   closed/manski_binary_xy
# MTR                    [0.550, 0.850]   fiber
# MTS                    [0.350, 0.700]   lp
# MTR & MTS              [0.550, 0.700]   lp
# Tan MSM, L=3           [0.467, 0.800]   closed/tan_msm_re1_L3
# Rosenbaum, G=9         [0.467, 0.800]   closed/anchored_rosenbaum_re1_G9
# c-dependence, c=0.1    [0.625, 0.750]   msm_agg/msm_transport_lp
# margin, kappa=0.05     [0.650, 0.750]   lp
# f-sens. KL, rho=0.1    [0.480, 0.846]   convex
# support k=3            infeasible       mip
\end{lstlisting}

\subsection{Ternary exposure}
\label{app:code-ternary}

A confounded treatment with three levels. Under the causal order alone the
natural bounds of \citet{manski1990nonparametric}, \eqref{eq:manski}, apply
for any finite $X$, with width $1-P(X{=}x)$.
\begin{lstlisting}
from partialid import Assumptions, Query, bound
from partialid.problems import ternary_exposure

p = ternary_exposure()            # X in {0,1,2}, Y in {0,1}
Q = Query.interventional(outcome="Y", do={"X": 1}, value=1)
a = Assumptions.order("X", "Y")   # natural bounds
r = bound(p.variables, p.data, Q, a)
print(f"[{r.lower:.3f}, {r.upper:.3f}] width={r.width:.3f} tight={r.tight} method={r.method}")
print(f"1 - P(X=1) = {p.expected_width:.3f}")
# [0.158, 0.730] width=0.573 tight=True method=fiber
# 1 - P(X=1) = 0.573
\end{lstlisting}

\subsection{Higher cardinality: four treatment levels}
\label{app:code-card42}

The same natural bound \citep{manski1990nonparametric}, now with $X\in\{0,1,2,3\}$.
\begin{lstlisting}
from partialid import Assumptions, Query, bound
from partialid.problems import ternary_cardinalities

p = ternary_cardinalities(card_x=4, card_y=2)   # X in {0,..,3}, Y in {0,1}
Q = Query.interventional(outcome="Y", do={"X": 1}, value=1)
r = bound(p.variables, p.data, Q, Assumptions.order("X", "Y"))
print(f"[{r.lower:.3f}, {r.upper:.3f}] width={r.width:.3f} tight={r.tight} method={r.method}")
print(f"1 - P(X=1) = {p.expected_width:.3f}")
# [0.023, 0.882] width=0.859 tight=True method=fiber
# 1 - P(X=1) = 0.859
\end{lstlisting}

\subsection{Higher cardinality: ternary treatment and outcome}
\label{app:code-card33}

The query targets one level of a ternary outcome. The same natural bound
of \citet{manski1990nonparametric} holds.
\begin{lstlisting}
from partialid import Assumptions, Query, bound
from partialid.problems import ternary_cardinalities

p = ternary_cardinalities(card_x=3, card_y=3)   # X, Y in {0,1,2}
Q = Query.interventional(outcome="Y", do={"X": 1}, value=1)   # P(Y_{x=1} = 1)
r = bound(p.variables, p.data, Q, Assumptions.order("X", "Y"))
print(f"[{r.lower:.3f}, {r.upper:.3f}] width={r.width:.3f} tight={r.tight} method={r.method}")
print(f"1 - P(X=1) = {p.expected_width:.3f}")
# [0.187, 0.861] width=0.674 tight=True method=fiber
# 1 - P(X=1) = 0.674
\end{lstlisting}

\subsection{Binary instrument}
\label{app:code-iv}

The encouragement design of \cref{sec:ill-designs,ex:iv-lp}, the binary
instrument of \citet{balke1997bounds}. The classifier reports
degree one after the directed cut $I=\{Z\}$, $R=\{X,Y\}$. Dropping the
instrument assumptions (order only) widens the interval.
\begin{lstlisting}
from partialid import Assumptions, Query, bound, classify_program
from partialid.problems import binary_iv

p = binary_iv(seed=2)
Q = Query.interventional(outcome="Y", do={"X": 1}, value=1)
a = Assumptions.iv("Z", "X", "Y")          # Z -> X -> Y, X <-> Y
r = bound(p.variables, p.data, Q, a)
print(f"[{r.lower:.3f}, {r.upper:.3f}] width={r.width:.3f} tight={r.tight} method={r.method}")
d = classify_program(a.diagram(), query=Q)
print(d.kind.value, d.degree, d.wl_wr_cut)

# the same table under order only (one latent on Z, X, Y)
r0 = bound(p.variables, p.data, Q, Assumptions.order("Z", "X", "Y"))
print(f"order only: [{r0.lower:.3f}, {r0.upper:.3f}]")
# [0.283, 0.615] width=0.333 tight=True method=closed/balke_pearl_iv
# linear 1 (('Z',), ('X', 'Y'))
# order only: [0.154, 0.705]
\end{lstlisting}

\subsection{Two instruments}
\label{app:code-two-iv}

Two independent instruments violate Condition~3 of
\citet[Proposition~2]{symbolicbounds2023general} but still admit the directed
cut $I=\{Z_1,Z_2\}$ (\cref{sec:sachs}).
\begin{lstlisting}
from partialid import Assumptions, Query, bound, classify_program
from partialid.problems import two_instruments

p = two_instruments()
Q = Query.interventional(outcome="Y", do={"X": 1}, value=1)
a = Assumptions.diagram(
    ["Z1", "Z2", "X", "Y"],
    directed=[("Z1", "X"), ("Z2", "X"), ("X", "Y")],
    bidirected=[("X", "Y")],
)
r = bound(p.variables, p.data, Q, a)
print(f"[{r.lower:.3f}, {r.upper:.3f}] width={r.width:.3f} tight={r.tight} method={r.method}")
d = classify_program(a.diagram(), query=Q)
print(d.kind.value, d.degree, d.wl_wr_cut)
# [0.397, 0.789] width=0.392 tight=True method=lp
# linear 1 (('Z1', 'Z2'), ('X', 'Y'))
\end{lstlisting}

\subsection{Two treatments under one latent}
\label{app:code-multi}

A multi-cause design of the kind treated by \citet{shridharan2023scalable}: two
treatments $T_0\to T_1$, both causing $Y$, with one latent shared by all three.
One hyperedge covers the ancestral set, so covering holds with a single hyperedge (\cref{thm:lp}).
\begin{lstlisting}
from partialid import Assumptions, Query, bound
from partialid.problems import shridharan_multi_treatment

p = shridharan_multi_treatment(n_treatments=2, card=2)   # T0, T1, Y
Q = Query.interventional(outcome="Y", do={"T0": 1}, value=1)
a = Assumptions.diagram(
    ["T0", "T1", "Y"],
    directed=[("T0", "Y"), ("T1", "Y"), ("T0", "T1")],
    bidirected=[("T0", "T1", "Y")],         # one latent on all three
)
r = bound(p.variables, p.data, Q, a)
print(f"[{r.lower:.3f}, {r.upper:.3f}] width={r.width:.3f} tight={r.tight} method={r.method}")
# [0.367, 0.794] width=0.428 tight=True method=fiber
\end{lstlisting}

\subsection{Order-only chains}
\label{app:code-chains}

Under a causal order $X\prec Z_0\prec\cdots\prec Y$ with maximal confounding,
the mediators are inert for $P(Y_x)$. The bound is the natural bound of
\citet{manski1990nonparametric} on the $(X,Y)$-margin, and the package prunes
the mediators before solving.
\begin{lstlisting}
from partialid import Assumptions, Query, bound
from partialid.problems import chain_order_only

p = chain_order_only(3, card=2)          # X, Z0, Y binary
Q = Query.interventional(outcome="Y", do={"X": 1}, value=1)
r = bound(p.variables, p.data, Q, Assumptions.order("X", "Z0", "Y"))
print(f"[{r.lower:.3f}, {r.upper:.3f}] width={r.width:.3f} tight={r.tight} method={r.method}")
print(f"1 - P(X=1) = {p.expected_width:.3f}")
# [0.445, 0.820] width=0.374 tight=True method=fiber
# 1 - P(X=1) = 0.374
\end{lstlisting}
A longer binary chain and a ternary chain change only the constructor and the
order:
\begin{lstlisting}
p = chain_order_only(4, card=2)          # X, Z0, Z1, Y binary
r = bound(p.variables, p.data, Q, Assumptions.order("X", "Z0", "Z1", "Y"))
# [0.241, 0.761] width=0.520 tight=True method=fiber   (1 - P(X=1) = 0.520)

p = chain_order_only(3, card=3)          # X, Z0, Y in {0,1,2}
r = bound(p.variables, p.data, Q, Assumptions.order("X", "Z0", "Y"))
# [0.095, 0.722] width=0.627 tight=True method=fiber   (1 - P(X=1) = 0.627)
\end{lstlisting}

\subsection{Front door}
\label{app:code-fd}

The front-door criterion of \citet{pearl1995diagrams}: $X$ affects $Y$ only
through $Z$, and the latent is shared by $X$ and $Y$. The functional is a point.
\begin{lstlisting}
from partialid import Assumptions, Query, Variables, bound

V = Variables.from_binary("X", "Z", "Y")
pz1, py1 = {0: 0.4, 1: 0.7}, {0: 0.3, 1: 0.8}      # P(Z=1|x), P(Y=1|z)
P = {(x, z, y): 0.5 * (pz1[x] if z else 1 - pz1[x]) * (py1[z] if y else 1 - py1[z])
     for x in (0, 1) for z in (0, 1) for y in (0, 1)}
Q = Query.interventional(outcome="Y", do={"X": 1}, value=1)
fd = Assumptions.diagram(["X", "Z", "Y"], directed=[("X", "Z"), ("Z", "Y")],
                         bidirected=[("X", "Y")])
r = bound(V, P, Q, fd)
print(f"[{r.lower:.3f}, {r.upper:.3f}] tight={r.tight} method={r.method}")
r0 = bound(V, P, Q, Assumptions.order("X", "Z", "Y"))
print(f"order only: [{r0.lower:.3f}, {r0.upper:.3f}]")
# [0.650, 0.650] tight=True method=id
# order only: [0.325, 0.825]
\end{lstlisting}

\subsection{Measurement error}
\label{app:code-me}

\begin{lstlisting}
from partialid import Assumptions, Query, Variables, bound

V = Variables.from_binary("X", "Y", "Y2", observed=["X", "Y2"])   # Y is latent
P = {(0, 0): 0.30, (0, 1): 0.20, (1, 0): 0.15, (1, 1): 0.35}      # table of (X, Y2)
Q = Query.ate(outcome="Y", treatment="X", y=1, t1=1, t0=0)
a = (Assumptions.diagram(["X", "Y", "Y2"], directed=[("X", "Y"), ("Y", "Y2")],
                         bidirected=[("Y", "Y2")])
     & Assumptions.mtr("Y2", "Y"))
r = bound(V, P, Q, a)
print(f"[{r.lower:.3f}, {r.upper:.3f}] tight={r.tight} method={r.method}")
# [-0.400, 1.000] tight=True method=lp
\end{lstlisting}

\subsection{The multi-site panel}
\label{app:code-panel}

\begin{lstlisting}
import pandas as pd
from partialid import Assumptions, Query, bound, bound_ci, from_dataframe

cols = ["site", "age_bin", "educ", "prior_emp", "urban", "spouse", "treated", "employed"]
df = pd.read_csv("notebooks/data/job_training_panel.csv")        # n = 5000
order = Assumptions.order(*cols)
Q = Query.interventional(outcome="employed", do={"treated": 1}, value=1)

d = from_dataframe(df, columns=cols)
r = bound(d.variables, d.P, Q, order)
print(f"pooled:   [{r.lower:.3f}, {r.upper:.3f}]")

Qu = Query.interventional(outcome="employed", do={"treated": 1}, value=1,
                          condition={"urban": 1})
r = bound(d.variables, d.P, Qu, order)
print(f"urban=1:  [{r.lower:.3f}, {r.upper:.3f}]")

s0 = df[df["site"] == 0]                                          # n = 1640
d0 = from_dataframe(s0, columns=cols)
a = order & Assumptions.mtr("employed", "treated")
r = bound(d0.variables, d0.P, Q, a)
ci = bound_ci(d0.variables, d0.P, Q, a, n=len(s0), n_boot=200, seed=0)
print(f"site 0, MTR: [{r.lower:.3f}, {r.upper:.3f}]  95% CI [{ci.lower:.3f}, {ci.upper:.3f}]")
# pooled:   [0.369, 0.862]
# urban=1:  [0.403, 0.860]
# site 0, MTR: [0.641, 0.866]  95% CI [0.622, 0.881]
\end{lstlisting}

\subsection{The directed chain under four latent structures}
\label{app:code-chain}

\begin{table}[ht]
\centering
\caption{Directed chain $X\to W\to Y$ on~\eqref{eq:re3}. The classification
is that of \texttt{classify\_program}, with the program degree in
parentheses.}
\label{tab:audit}
\small
\begin{tabular}{@{}lllll@{}}
\toprule
Latents & Query & Classification & Interval & \texttt{tight} \\
\midrule
$\{X,W\}$ & $P(Y_1{=}1)$ & LP (1) & $[0.470,0.570]$ & yes \\
$\{X,W\}$ & PNS & bilinear (2) & $[0.000,0.520]$ & no \\
$\{X,W\},\{W,Y\}$ & $P(Y_1{=}1)$ & LP (2) & $[0.270,0.770]$ & yes \\
$\{X,W\},\{W,Y\}$ & PNS & LP (2) & $[0.000,0.520]$ & yes \\
$\{X,W,Y\}$ & $P(Y_1{=}1)$ & LP (1) & $[0.270,0.770]$ & yes \\
three pairwise edges & $P(Y_1{=}1)$ & polynomial (3) & $[0.270,0.770]$ & no \\
\bottomrule
\end{tabular}
\end{table}

The rows of \cref{tab:audit}. The mediated PNS and the triangle are not
certified, and the code prints \texttt{tight=False} for them.
\begin{lstlisting}
from partialid import Assumptions, Query, Variables, bound, classify_program

V = Variables.from_binary("X", "W", "Y")
Pxw = {(0, 0): 0.25, (0, 1): 0.25, (1, 0): 0.15, (1, 1): 0.35}
pyw = {0: 0.4, 1: 0.6}                                            # P(Y=1|w)
P = {(x, w, y): Pxw[x, w] * (pyw[w] if y else 1 - pyw[w])
     for x in (0, 1) for w in (0, 1) for y in (0, 1)}
queries = {"P(Y_1=1)": Query.interventional(outcome="Y", do={"X": 1}, value=1),
           "PNS": Query.pns(treatment="X", outcome="Y")}
latents = {"mediated": [("X", "W")],
           "pair chain": [("X", "W"), ("W", "Y")],
           "one oval": [("X", "W", "Y")],
           "triangle": [("X", "W"), ("W", "Y"), ("X", "Y")]}
for name, bi in latents.items():
    a = Assumptions.diagram(["X", "W", "Y"], directed=[("X", "W"), ("W", "Y")],
                            bidirected=bi)
    for qn, Q in queries.items():
        k = classify_program(a.diagram(), query=Q, variables=V)
        r = bound(V, P, Q, a)
        print(f"{name:10s} {qn:8s} {k.kind.value:15s} {k.degree} "
              f"[{r.lower:.3f}, {r.upper:.3f}] tight={r.tight}")
# mediated   P(Y_1=1) linear          1 [0.470, 0.570] tight=True
# mediated   PNS      bilinear/factor 2 [0.000, 0.520] tight=False
# pair chain P(Y_1=1) linear          2 [0.270, 0.770] tight=True
# pair chain PNS      linear          2 [0.000, 0.520] tight=True
# one oval   P(Y_1=1) linear          1 [0.270, 0.770] tight=True
# one oval   PNS      linear          1 [0.000, 0.520] tight=True
# triangle   P(Y_1=1) polynomial      3 [0.270, 0.770] tight=False
# triangle   PNS      polynomial      3 [0.000, 0.520] tight=False
\end{lstlisting}

\clearpage
